\input{../tex-inputs/latex-leseansicht-vorspann}

\section[Gustav Schwarzkopf an Arthur Schnitzler, 12. 1. 1897]{L04285 Gustav Schwarzkopf an Arthur Schnitzler, 12. 1. 1897}
\nopagebreak\mylabel{L04285v}
\rehead{ }\normalsize\beginnumbering\briefempfaengerindex{Schnitzler, Arthur@\textsc{Schnitzler, Arthur}!zzzSchwarzkopf, Gustav@\emph{von Gustav Schwarzkopf}!1897-01-121@{12. 1. 1897}|(be}
\toendnotes[C]{\smallbreak\pagebreak[2]}
\correspDesc{Versand  durch Gustav Schwarzkopf am 12. 1. 1897 in Wien
\newline{}Übermittlung  am 12. 1. 1897 in Wien
\newline{}Zustellung  am 12. 1. 1897 in Wien
\newline{}Erhalt  durch Arthur Schnitzler am 12. 1. 1897 in Wien}\toendnotes[C]{\smallbreak}
\Standort{CUL, Schnitzler, B 96a.}
\physDesc{Postkarte, 193 Zeichen
\newline{}Handschrift: schwarze Tinte, deutsche Kurrent
\newline{}Versand: 1) Stempel: »\nobreak{}\oindex{IX., Alsergrund@\textbf{IX., Alsergrund}, \emph{Verwaltungsgebiet}|pwk}Wien \textcolor{gray}{9}/1 66, 12. 1. 97, 1–2N\nobreak{}«.   2) Stempel: »\nobreak{}\oindex{IX., Alsergrund@\textbf{IX., Alsergrund}, \emph{Verwaltungsgebiet}|pwk}Wien 9/3 72, 12. 1. 97, 3.N, Bestellt\nobreak{}«. }\pstart{}{\pb}Herrn \textsc{D\textsuperscript{r} Arthur Schnitzler}\pend{}\pstart{}\textsc{Wien\oindex{Wien@\textbf{Wien}, \emph{Verwaltungsgebiet}|pw}}\pend{}\pstart{}IX Frankgaſſe 1\oindex{Wien@\textbf{Wien}!IX., Alsergrund@\textbf{IX., Alsergrund}!Frankgasse 1@\textbf{Frankgasse 1}, \emph{Wohngebäude}|pw}.\pend{}{\bigskip}\vspace{1em}
\pstart
           \raggedleft{}{\pb}12. I. 97\pend
           \vspace{0.5em}
\pstart
           Lieber Arthur, es iſt mir leider nicht möglich zu ko{\geminationm}en, ich bin mit \textsc{\uline{Ganghofer}’s\pwindex{Ganghofer, Katharina 7.\,7.\,1856 Pest – 8.\,4.\,1930 München@\textsc{Ganghofer, Katharina} (7.\,7.\,1856 Pest – 8.\,4.\,1930 München), \emph{Schauspielerin}|pw}\pwindex{Ganghofer, Ludwig 7.\,7.\,1855 Kaufbeuren – 24.\,7.\,1920 Tegernsee@\textsc{Ganghofer, Ludwig} (7.\,7.\,1855 Kaufbeuren – 24.\,7.\,1920 Tegernsee), \emph{Schriftsteller}|pw}}, die morgen abreiſen, bei \textsc{Weiss\pwindex{Karlweis, Carl 23.\,11.\,1850 Wien – 27.\,10.\,1901 ebd.@\textsc{Karlweis, Carl} (23.\,11.\,1850 Wien – 27.\,10.\,1901 ebd.), \emph{Schriftsteller}|pw}} geladen.\pend
           
\pstart
           Herzlichſt Ihr{\\[\baselineskip]}\spacefill\mbox{GustavSch}\pend
           \leftskip=0em{}\selectlanguage{ngerman}\endnumbering\briefempfaengerindex{Schnitzler, Arthur@\textsc{Schnitzler, Arthur}!zzzSchwarzkopf, Gustav@\emph{von Gustav Schwarzkopf}!1897-01-121@{12. 1. 1897}|)be}\mylabel{L04285h}
\begin{anhang}
\end{anhang}\newcommand{\dateiname}{L04285}\newcommand{\titel}{Gustav Schwarzkopf an Arthur Schnitzler, 12. 1. 1897}\newcommand{\editorInnen}{Herausgegeben von Jahnke, SelmaMüller, Martin Anton}\input{../tex-inputs/latex-leseansicht-abspann}
